\subsection{Necessary structure of the remaining quadratic pair}
\label{sec:m5-remaining-quadratic}
The rings $R7\_5\_3895$ and $R7\_5\_5039$ remain unresolved at
this checkpoint.  Proposition~\ref{prop:m5-quadratic-galois-pu} reduces
any categorification to a pseudo-unitary Galois conjugate.  The following
restrictions apply to that positive representative.  They do not by
themselves give a categorification or a general exclusion.  The final
proposition excludes braided categorifications of both rings.

\paragraph{Full character transport.}
Write $d_i=a_i+b_i\sqrt5$.  In source order the vectors are
\[
\begin{array}{c|l|l}
 & (a_i) & (b_i)\\\hline
3895&(1,5,5,10,8,8,9)&(0,2,2,4,4,4,4)\\
5039&(1,5,5,8,8,10,9)&(0,2,2,4,4,4,4).
\end{array}
\]
Both global dimensions are
$D=720+320\sqrt5=80(9+4\sqrt5)$.
Let $P$ swap labels $4,5$ for the first ring and labels $3,4$ for
the second.  The signed normalized Galois transport of a column restriction
vector is the integral involution
\begin{equation}\label{eq:m5-q5-transport}
 T=P+\mathbf b\left(\mathbf a-\frac94\mathbf b\right)^{\!t}.
\end{equation}
The permutation term records the action on the non-Frobenius--Perron
characters; omitting it gives an incorrect transport.  The exact checker
constructs the Galois automorphism on the separating fusion matrix by
Chinese remaindering over its character factors.  If $A$ is this algebra
automorphism and $B$ is the Casimir, the $i$th column of $T$ is the
transpose of row zero of $B A(B)^{-1}A(N_i)$.  Evaluating in all source
characters verifies
\[
 CT=\operatorname{diag}(f)P_\sigma\operatorname{diag}(1/f)C.
\]
This checks the complete transport without numerical root choices or
an assumed labelling of repeated formal codegrees.  In particular
$Te_0=(1,2,2,4,4,4,4)^t$, the forced restriction row of the Galois
image of the central unit.  Complete row transport, character reciprocity
and duality initially leave two low induction frames for each source ring.
The refinements below remove one frame in each case.
For a positive restriction row, its actual transported row is the
nonnegative choice among $Tv$ and $-Tv$.  Orbit counts must retain
this sign; the linear map alone need not preserve nonnegativity.

\paragraph{A second saturated quadratic character.}
Besides the Frobenius--Perron pair, each source has a pair of formal
codegree-five characters exchanged by $\sqrt5\mapsto-\sqrt5$.
In the orthonormal basis of the two corresponding character-difference
vectors, the full known induced-unit compression is
\[
 B_5=\begin{pmatrix}-\sqrt5/10&0\\0&\delta\sqrt5/5\end{pmatrix},
 \qquad \frac15I-B_5^*B_5=
 \begin{pmatrix}3/20&0\\0&0\end{pmatrix},\qquad \delta=\pm1.
\]
Thus the codegree-five difference saturates the conductor-five norm
bound even though the Frobenius--Perron difference does not.  The sign
$\delta$ is determined by each exact character assignment.

The equality case gives a criterion that is independent of these two
rings and of the prime five.

\begin{proposition}[Saturated quadratic compression fixes primitive phases]
\label{prop:prime-quadratic-saturation}
Let $p\equiv1\pmod4$ be prime.  In an anomaly-one unitary congruence
representation of $\mathrm{SL}_2(\mathbb Z)$, let $w$ be fixed by the
upper unipotents and have diagonal-torus character the primitive
quadratic character $\chi_p$ of conductor $p$.  Let $L$ be a subspace
of the twist-one space containing $w$.  Suppose
\[
 \operatorname{proj}_L(Sw)=\frac{\delta}{\sqrt p}w,
 \qquad \delta\in\{1,-1\}.
\]
Then the cyclic representation generated by $w$ factors through
$\mathrm{SL}_2(\mathbb F_p)$, and the full twist-one projection of
$Sw$ equals the displayed known projection.  Every nontrivial twist
in the support of $Sw$ is a primitive $p$th root $\zeta_p^k$ with
$\chi_p(k)=\delta$.
\end{proposition}
\begin{proof}
We first determine the equality space at every modulus $p^a$.
The upper-unipotent subgroup $U$ is the stabilizer of $(1,0)$ in
the transitive action on primitive column vectors $(r,s)$ modulo
$p^a$.  Frobenius reciprocity embeds every irreducible with a nonzero
$U$-fixed vector isometrically in the corresponding permutation
representation.  In its $U$-fixed, $\chi_p$-torus-character subspace,
the functions supported on $s$ a unit form the line generated by
\[
 F(r,s)=\chi_p(s)[p\nmid s].
\]
Write $\Sigma$ for its orthogonal complement, supported on the small
cell $p\mid s$, where $r$ is a unit.  Let $P_U$ be unipotent averaging
and $A=P_USP_U$ the Weyl compression.  Since $\chi_p(-1)=1$, this
compression is self-adjoint.  Direct averaging gives
\[
 AF=G,\qquad G(r,s)=\chi_p(-r)[p\mid s].
\]
Indeed, $\chi_p(r+ts)$ is constant if $p\mid s$ and averages to
zero otherwise.  The Weyl translate of a function in $\Sigma$
vanishes on the small cell after averaging, so $A\Sigma\subseteq
\mathbb CF$.  Moreover $\|G\|^2/\|F\|^2=1/p$.
Self-adjointness now gives $AG=F/p$ and makes $A$ vanish on
$\Sigma\cap G^\perp$.  Consequently $\|A\|=1/\sqrt p$, and its
entire maximal-singular subspace is $\operatorname{span}(F,G)$.
Both functions pull back from level $p$, as does the invariant
subrepresentation they generate.  Equality therefore forces level
$p$ on the cyclic subrepresentation of the equality vector.

At another prime the local torus character is trivial and the
compression bound is one.  Equality says that both the vector and
its Weyl translate are upper-unipotent-fixed; hence the vector is
fixed by upper and lower unipotents.  These generate the local
$\mathrm{SL}_2$, so this factor is trivial.  Chinese remaindering and
orthogonal irreducible decomposition give the global norm bound
$1/\sqrt p$ and the same equality conclusion: the top eigenspace of
a tensor product of positive operators $A^*A$ is the tensor product
of their top eigenspaces.  Finally, the hypothesized projection onto
$L$ already has squared norm $\|w\|^2/p$.  Since $L$ lies in the
full twist-one space, Pythagoras forces all additional twist-one
components to vanish.  This is an equality refinement of the local
compression calculation in \cite[Lemma~4.1]{SchopierayDimensionBlock}.

Put $h=\|w\|^2$.  The twist-one squared weight of $Sw$ is $h/p$;
let $x_+,x_-$ be its remaining squared weights on the two quadratic
classes of primitive $p$th roots.  Thus
$x_++x_-=h(p-1)/p$.  The diagonal torus acts transitively within
each class, so the squared weights there are equal.  The average
phases in the two classes are $(-1\pm\sqrt p)/(p-1)$, by the
quadratic Gauss sum.  Also $S^{-1}w=Sw$, because $S^2$ acts by
$\chi_p(-1)=1$.  Using $Tw=w$ and
$STS=T^{-1}ST^{-1}$ gives
\[
 \langle Sw,TSw\rangle=\langle w,Sw\rangle
     =\frac{\delta h}{\sqrt p}.
\]
Consequently
\[
 \frac hp+x_+\frac{-1+\sqrt p}{p-1}
            +x_-\frac{-1-\sqrt p}{p-1}
       =\frac{\delta h}{\sqrt p},
 \qquad x_+-x_-=\frac{\delta h(p-1)}p.
\]
Nonnegativity implies $x_{-\delta}=0$, proving the phase assertion.
\end{proof}

For the two present rings apply the proposition to the unnormalized
codegree-five difference $w$, of squared norm two.  Its known
twist-one projection has squared norm $2/5$.  Thus every additional
center row with unequal coordinates $4,5$ in the first source, or
$3,4$ in the second, has primitive fifth-root twist of residue class
$\delta$.  In fact $(Sw)_Y=(v_p-v_q)/\sqrt5$, up to the sign of
$w$, for the indicated pair $(p,q)$.  Unitarity gives the small
integral budget
\[
 \sum_{Y\text{ additional}}(v_p-v_q)^2=8.
\]
Each permitted primitive eigenvalue contributes four.  Hence every
additional row satisfies $|v_p-v_q|\leq2$; a transported orbit with
difference two exhausts this budget.

There is one further cut for $R7\_5\_5039$.  Its characters
\[
 (1,-1\pm\sqrt{-2},-1\mp\sqrt{-2},0,0,0,1)
\]
have formal codegree eight.  Their difference has odd conductor-eight
torus character.  If $v$ is the restriction row of one of the two
induced-unit constituents, the compression is
$\sqrt{-2}(v_1-v_2)/4$, of absolute value at most $1/\sqrt8$.
Thus $|v_1-v_2|\leq1$.  One of the two frames violates this inequality;
the other, with three identical codegree-eight rows, survives it.
This uses the local bound for all torus characters, including the odd
quadratic character, and does not follow merely from the real quadratic
Casimir criterion.

\paragraph{A conductor-eight partial-projection criterion.}
The odd character above also gives a uniform refinement that records
small two-primary twist orders, including twists with arbitrary odd
parts.  The distinction between these two projectors is essential.

\begin{proposition}[Odd conductor-eight projection ladder]
\label{prop:chi-eight-projection-ladder}
Let $\chi_{-8}$ be the primitive quadratic character modulo eight,
with values $(1,1,-1,-1)$ at $(1,3,5,7)$.  Let $w$ be a
unipotent-fixed vector of this torus character in a unitary congruence
representation.  For $r=0,1,2,3$ let $Q_r$ project onto invariants of
the two-primary unipotent subgroup generated by $T^{2^r}$, acting as
the identity on the odd-primary factors.  If $b=\|Q_0Sw\|^2$, then
\[
 \|Q_rSw\|^2=2^rb,\qquad b\leq\frac{\|w\|^2}{8}.
\]
Thus the squared masses of $Sw$ with two-primary twist order exactly
$1,2,4,8$ are respectively $b,b,2b,4b$.  The same upper bounds hold if
$Q_r$ is replaced by the full projector onto $T^{2^r}$-invariants.
Equality in a two-primary bound forces that local cyclic factor to
descend to level eight.  Equality in the full bound also makes all
odd-primary cyclic factors trivial.
\end{proposition}
\begin{proof}
At modulus $M=2^a$, $a\geq3$, use the primitive-vector universal
representation as in Proposition~\ref{prop:prime-quadratic-saturation}.
Its unipotent-fixed torus-character space is the orthogonal sum of
\[
 F(u,v)=\chi_{-8}(v)[v\text{ odd}]
\]
and the small-cell functions
$f(u,v)=\chi_{-8}(u)g(uv)$, supported on $v$ even and $u$ odd.
Unipotent invariance implies
\[
 g(z)=\chi_{-8}(1+tz)g(z+tz^2).
\]
The map $z\mapsto z+tz^2$ permutes the even residues and the
congruence classes below: its difference at $z,z'$ is
$(z-z')(1+t(z+z'))$, with odd second factor.  Taking $t=2$ on
$z\equiv2\pmod4$ and $t=1$ on $z\equiv4\pmod8$ gives
\[
 \sum_{z\equiv2\ (4)}g(z)=
 \sum_{z\equiv4\ (8)}g(z)=0,
\]
since the character multiplier is minus one in each case.
Use the Weyl convention $(Sf)(u,v)=f(-v,u)$ and write
$R_r=P_{\langle T^{2^r}\rangle}S$ locally.  Direct averaging gives
\[
 R_rF(u,v)=\chi_{-8}(u)[v\equiv0\pmod{2^{3-r}}],
\]
and, on $v$ odd and $u\equiv0\pmod{2^r}$,
\[
 R_rf(u,v)=\chi_{-8}(-v)\frac{2^r}{M}
                 \sum_{z\equiv0\ (8)}g(z),
\]
with value zero elsewhere.  The two cancellation identities justify
retaining only the residues divisible by eight in the latter sum.
The first output has squared norm $2^{r-3}\|F\|^2$.  The second
has $M^2/2^{r+1}$ possible nonzero positions, and therefore
\[
 \|R_rf\|^2=2^{r-1}\left|\sum_{z\equiv0\ (8)}g(z)\right|^2
 \leq M2^{r-4}\sum_{z\text{ even}}|g(z)|^2
 =2^{r-3}\|f\|^2.
\]
For $r=0,1,2$, the outputs of the two cells have opposite parity of
$v$, so are orthogonal.  Their exact norm expressions give
$R_r^*R_r=2^rR_0^*R_0$.  To include $r=3$, use $t=2$ separately
on $z\equiv2,6\pmod8$; each class is preserved and each sum
vanishes.  Only the zero class modulo eight survives.  The small-cell
formula above now has squared norm $4|\sum_{8\mid z}g(z)|^2$.
The big-cell image is $\chi_{-8}(u)[u\text{ odd}]$ with squared
norm $\|F\|^2$.  These images are disjoint by parity of $u$,
and both squared norms are eight times their $r=0$ values.  Thus
$R_3^*R_3=8R_0^*R_0$ as well.
Equality on the small cell requires $g$ to be constant on the
multiples of eight and zero elsewhere.  Thus the complete extremal
space is spanned by $F$ and $\chi_{-8}(u)[8\mid v]$, both pulled
back from level eight.

Frobenius reciprocity and orthogonal decomposition transfer these
identities to arbitrary local unitary representations.  For $Q_r$,
the odd-primary Weyl operators are unitary, preserving both the exact
ratios and the bound.  For the full projector they are instead
unipotent compressions of norm at most one.  Equality there forces
upper- and lower-unipotent invariance at each odd prime.  The asserted
factorization and triviality conclusions follow as in
Proposition~\ref{prop:prime-quadratic-saturation}.
\end{proof}

For $R7\_5\_5039$, the character-difference vector $w_8$ has squared
norm two and
$(Sw_8)_Y=\sqrt{-2}(v_1-v_2)/4$.  Hence
\[
 \sum_Y(v_1-v_2)^2=16.
\]
In squared-difference units, the exact two-primary masses at orders
$1,2,4,8$ are $c,c,2c,4c$ with $0\leq c\leq2$.  Once the exponent prime support
$\{2,5\}$ is imposed below, the combined nonzero phase groups of
root orders $\{1,5\}$, $\{2,10\}$ and $\{4,20\}$ consequently
have budgets at most $2,2,4$.  They are inequalities: a root whose
odd part is divisible by $25$ has zero first and second relative
trace and can carry additional mass in the same two-primary class.
The independent audit and finite convention checks are in
\path{evidence/rank7_m5_20261020/modules/CHI_MINUS8_PARTIAL_PROJECTION_AUDIT.md}.

\paragraph{Exact phase values and a forced third rotation.}
By the Cauchy theorem for spherical categories
\cite[Theorem~3.9]{BNRW}, the Frobenius--Schur exponent has only the
prime divisors $2,5$, since $D/80$ is an algebraic unit.  Thus every
central twist order is $2^a5^b$.  For normalized relative trace $t$ to
$\mathbb Q(\sqrt5)$, the pair $(t(\theta),t(\theta^2))$ has exactly
the following ten possible values.  In coefficient notation
$(r,s,t,u)=(r+s\sqrt5,t+u\sqrt5)$ they are
\begin{gather*}
 (1,0,1,0),\quad(-1,0,1,0),\quad(0,0,-1,0),\quad(0,0,0,0),\\
 P_\pm=(-1/4,\pm1/4,-1/4,\mp1/4),\\
 Q_\pm=(1/4,\pm1/4,-1/4,\pm1/4),\qquad
 R_\pm=(0,0,1/4,\pm1/4).
\end{gather*}
Nonzero pairs come from orders $1,2,4,5,10,20$; larger allowed
prime powers give zero in both coordinates by the Ramanujan or induced
quadratic Gauss-sum formula.  These are exact trace values, not only
vertices of a containing convex hull.  Counts at each value on an
orbit of restriction rows are therefore nonnegative integers.

Since three is prime to the exponent, the third relative twist moment
is the field conjugate of the first.  Put $h=D/\bar D$.  The complete
row transport implies that each dimension-weighted orbit sum $W$
satisfies $W=hT\bar W$: the row and its dimension acquire the same
transport sign, which cancels in their product.  The full first induction moment is
$L_1=De_0$, including the induced-unit constituents.  Hence
\[
 L_3=hT\overline{L_1}=DTe_0,
 \qquad (\nu_3(X_i))_i=(1,2,2,4,4,4,4).
\]
The trace identity is \cite[Theorem~4.1]{NgSchauenburgExponent}.
Each entry equals $\dim\operatorname{Hom}(1,X_i^3)$, so every third
rotation operator is the identity.  This is a necessary categorical
consequence and presently yields no contradiction.

\paragraph{A source-fixed conductor-eight dichotomy.}
For $R7\_5\_5039$, call a center simple active when
$k_Y=v_{Y,1}-v_{Y,2}\ne0$.  Source duality changes $k_Y$ to $-k_Y$,
while center duality preserves the twist.  Each active object therefore
has a distinct dual with the same squared weight.  The mass $c$ above
is consequently an even integer.  There are exactly two alternatives:
\[
\begin{array}{c|l|l}
 &\text{exact two-primary masses}&\text{possible active twist orders}\\\hline
c=2&2,2,4,8\text{ at orders }1,2,4,8&1,2,4,5,8,10,20,40\\
c=0&0\text{ at all orders at most }8&16,32,64,80,160,320.
\end{array}
\]
Here the finite order lists use Galois actions fixing the complete
source field, as follows; they do not assume that the rational
transport matrix has a positive sign on every row.

All source characters, codegrees and dimensions lie in
$L=\mathbb Q(\sqrt5,\sqrt{-2})$.  A cyclotomic automorphism fixing
$L$ fixes a center restriction row under the categorical Galois
permutation.  Indeed, its sign is positive from the fixed positive
entry $S_{0,Y}=\dim(Y)/D$, and the entries
$S_{A_E,Y}=\chi_E(v_Y)/f_E$ then determine $v_Y$ through the complete
invertible source character table.  The twist is transported by the
square of the automorphism.  For an object of twist order $2^a5^b$,
choose $u=3\pmod8$ at the two-primary factor and the identity at the
fifth-primary factor.  If $a\geq4$, then
$v_2(u^2-1)=3$ and this squared action has order $2^{a-3}$.
Independently, $u=6\pmod{25}$ at the fifth-primary factor, with the
identity at the two-primary factor, has squared action of order
$5^{b-1}$ when $b\geq2$.  Both choices fix $L$, and Chinese
remaindering makes them independent.  Thus the row-fixed orbit has
size divisible by
\[
 m(a,b)=2^{\max(a-3,0)}5^{\max(b-1,0)}.
\]
Its distinct source-dual row gives a disjoint orbit, so an active
object costs at least $2m(a,b)k_Y^2$ in its exact two-primary mass.
If $c=2$, the four masses exhaust the total sixteen.  Each is at
most eight, excluding $b\geq2$, which would cost at least ten.
If $c=0$, then $a\geq4$; $b\geq2$ would cost at least twenty,
and $a\geq7$ would cost at least thirty-two.  These contradictions
prove the displayed lists.  Both alternatives remain necessary data,
with neither presently excluded.

\paragraph{The full source phase domain.}
There is a stronger finite domain than the ten pairs over
$\mathbb Q(\sqrt5)$.  Put
$F=\mathbb Q(\sqrt5,\sqrt2,i)$ and $\xi=\zeta_8$.
For normalized relative trace to $F$, the complete first/second
moment domain consists of the following $37$ distinct pairs:
\[
\begin{gathered}
 (\xi^j c_1,\xi^{2j}c_2)\quad(0\leq j<8),\qquad
 (0,\xi^j c_1)\quad(j=1,3,5,7),\qquad (0,0),\\
 (c_1,c_2)\in\left\{(1,1),
 \left(\frac{-1+\delta\sqrt5}{4},
       \frac{-1-\delta\sqrt5}{4}\right):\delta=\pm1\right\}.
\end{gathered}
\]
The first $24$ pairs have root orders dividing $40$; the next $12$
have orders $16$ or $80$.  Larger allowed prime powers have zero
trace in both coordinates: the relevant cyclotomic kernels fixing
$F$ negate the first or second power at two-primary order at least
$32$, and average to zero at fifth-primary order at least $25$.
Exact averaging in $\mathbb Q[x]/(\Phi_{80}(x))$ verifies the finite
list and its distinctness.  Every squared cyclotomic automorphism
fixes $F$, so the relative phases are constant on transported row
orbits.

The same row-fixed actions prove additional integer restrictions.
Each of the twelve order-sixteen phase counts at each fixed row is
even.  The genuine zero-phase multiplicity at a fixed row belongs to
$4\mathbb Z_{\geq0}+5\mathbb Z_{\geq0}$: partition by twist order
and use the row-fixed cycles of length divisible by four or five.
For the active rows of $R7\_5\_5039$, the dichotomy sharpens these
conditions.  In the branch $c=2$ only the first $24$ phases are
allowed, with exact two-primary masses $2,2,4,8$.  In the branch
$c=0$ only the twelve order-sixteen phases and the zero phase are
allowed; the latter count at each active row is divisible by four,
since its possible orders are $32,64,160,320$.
The full all-order argument and exact field calculation are recorded
in \path{models/sqrt5/FULL_SOURCE_PHASES_PROOF.md} under the
continuation evidence directory.  The source-fixed dichotomy is
independently audited in \path{modules/CHI_MINUS8_SIGN_SAFE_DICHOTOMY.md}.

\paragraph{The distinct-row frame is excluded.}
The two possible low frames for $R7\_5\_3895$ differ only in the
restriction rows of its two codegree-sixteen constituents.  In the
first they are
\[
 (1,0,2,0,1,1,2),\qquad (1,2,0,0,1,1,2);
\]
in the second both are $(1,1,1,0,1,1,2)$.
The first frame is now excluded by a complete exact integer
certificate.  Here is its categorical and finite-domain content.

The phase multiplicities are nonnegative integers.  When the actual
positive quadratic Galois transport fixes a row or takes it to its source dual,
the six nonrational phase types over $\mathbb Q(\sqrt5)$ have even
multiplicity: compose the Galois permutation with duality if needed;
its odd iterates change the fifth-primary twist and cannot fix an
object of fifth-primary order five.  The relevant columns are
therefore doubled.  In addition to the restriction Gram matrix,
first and second indicator moments, and induced-unit $STS$
identities, impose all $56$ unordered triples among the six nonunit
induced-unit objects.  Their symmetric Verlinde coefficients are
\[
 \frac{D}{f_Ef_Hf_J}\sum_Y
 \frac{\chi_E(v_Y)\chi_H(v_Y)\chi_J(v_Y)}{\dim(Y)}
       \in\mathbb Z_{\geq0}.
\]
This is $N_{EH}^{J^*}$; the absence of a complex conjugate is
appropriate for the symmetric triple formulation.  Its irrational
coefficient vanishes and its rational coefficient is an integer,
encoded by a nonnegative integer variable.

Conductor-five saturation fixes the phase of every row with
$v_4\ne v_5$, and its squared-difference budget is eight.
After integral spectral filtering there are $58$ possible active
row orbits: $25$ have weight two, $27$ have weight eight, and six
have weight four with even multiplicity.  Hence there are exactly
\[
 \binom{28}{4}+27+6=20508
\]
active allocations.  Exact residual Gram bounds reduce these to
$126$ for the distinct-row frame.  Integer-lattice containment
reduces them to $32$ per positive-indicator character labelling.
There are four such labellings.  Exact Farkas covectors exclude
$28$ allocations in each, leaving $16$ combined branches.  Complete
integer split trees close all sixteen with $775$ terminal integer
Farkas inequalities.  At each node the split $x\leq k$ or
$x\geq k+1$ covers every integer value; at each terminal node the
covector is nonnegative on every allowed column and negative on
the residual right side.  Branch pruning uses explicit entry
bounds and negative quadratic-form witnesses for the residual
Gram matrix.  The negative-indicator cases are excluded by a
further exact covector.

The complete entry point
\path{evidence/rank7_m5_20261020/integral_gram/verify_frame6.py}
regenerates all $15$ initial low frames and $163766$ high rows from
the canonical source, both surviving orbit inventories, all eight
positive-indicator systems across the two frames, and every active
allocation.  It then replays the negative-indicator, lattice,
separation, pruning and integer-tree certificates.  Normal and
optimized Python both pass without an optimizer.  Its conclusion
is precisely that the distinct-row frame (index six) is impossible;
the repeated-row frame (index eight) remains open.  Thus every
hypothetical categorification must, after the established positive
Galois reduction, use the repeated codegree-sixteen restriction row.

The exact transport, low-block compression, and equality computations
are recorded in \path{evidence/rank7_m5_20261020/models/sqrt5/} and
the independent saturation audit in the adjacent \path{modules/}
directory.  Earlier integral phase and low-block Verlinde witnesses
certify their recorded truncated systems only; subsequent higher-word
constraints reject those particular witnesses.  They are not complete
modular data or models.  No coherent associators or general exclusion
have been established for either remaining tensor.

\begin{proposition}[The remaining pair has no braided categorification]
\label{prop:m5-frontier-no-braiding}
Neither $R7\_5\_3895$ nor $R7\_5\_5039$ is the Grothendieck ring of
a braided fusion category over an algebraically closed field of
characteristic zero.  Their categorification without a braiding remains
unresolved.
\end{proposition}
\begin{proof}
Proposition~\ref{prop:m5-quadratic-galois-pu} permits a pseudo-unitary
Galois conjugate, and a braiding is preserved under this operation.
In both source tensors each nonunit basis element generates the whole
ring.  Consequently the M\"uger center is either $\mathrm{Vec}$ or the
whole category.  The second possibility is symmetric and hence integral
\cite[Corollary~9.9.11]{EGNO}, contrary to the simple dimension
$5+2\sqrt5$.  The canonical positive spherical structure therefore
makes the category modular
\cite[Propositions~9.5.1 and~8.20.12]{EGNO}.

For $R7\_5\_5039$, modularity requires the formal codegree
$D/(5+2\sqrt5)^2=16$ \cite[Section~3]{OstrikCodegrees}.
Its exact Casimir polynomial is
\[
 (x-8)^3(x-5)^2(x^2-1440x+6400),
\]
which does not vanish at $16$.

For $R7\_5\_3895$, matching every simple dimension $d_j$ with the
formal codegree $D/d_j^2$ gives exactly four character assignments $p$.
Each forces a symmetric real orthogonal matrix
\[
 S_{ij}=\frac{\chi_{p(j)}(X_i)d_j}{20+8\sqrt5}.
\]
Its entries lie in $K=\mathbb Q(\sqrt5)$, and the nontrivial field
automorphism has character permutation $\pi=(0\ 6)(4\ 5)$ in
the zero-based source labelling.

We first bound all possible twist orders; no trial cutoff is used.
The categorical Cauchy theorem \cite[Theorem~3.9]{BNRW} permits only
the primes $2,5$, since $D=80(9+4\sqrt5)$ and $9+4\sqrt5$ is a unit.
For the actual twists, Galois symmetry gives
\begin{equation}\label{eq:m5-braided-galois-twist}
 \frac{\theta_{\pi_\sigma(i)}}{\theta_{\pi_\sigma(0)}}
       =\sigma^2(\theta_i)
\end{equation}
by \cite[Proposition~4.7, equation~(4.12)]{DongLinNg}.
The denominator is retained; the modular category need not have
trivial anomaly.  An automorphism fixing $K$ has $\pi_\sigma=1$
and therefore square-fixes every twist.  Taking exponent $3$ on the
two-primary part and $1$ on the fifth-primary part gives
$v_2(3^2-1)=3$.  Taking exponent $6$ on the fifth-primary part and
$1$ on the two-primary part gives $v_5(6^2-1)=1$.  Both choices fix
$K$, so every twist has order dividing $40$.

The exponent $17$ modulo $40$ acts nontrivially on $K$ and squares
to $9$.  Writing $\theta_i=\zeta_{40}^{k_i}$, equation
\eqref{eq:m5-braided-galois-twist} becomes
$k_{\pi(i)}=k_6+9k_i\pmod{40}$, with $k_0=0$.
Its complete solution set is
\[
 \begin{gathered}
 k_6=8c,\quad 0\leq c<5,\qquad
 k_1,k_2,k_3\equiv-c\pmod5,\\
 0\leq k_4<40,\qquad k_5\equiv8c+9k_4\pmod{40}.
 \end{gathered}
\]
There are $5\cdot8^3\cdot40=102400$ candidates per matrix.
Exact reduction modulo $\Phi_{40}$ imposes every balancing identity
\[
 \theta_i\theta_j\sqrt D\,S_{ij}
       =\sum_kN_{ij}^k d_k\theta_k.
\]
Only four assignments remain: with $\zeta=\zeta_5$, their twists are
\[
 (1,-1,-1,1,\zeta^2,\zeta^3,1),\qquad
 (1,-1,-1,-1,\zeta,\zeta^4,1),
\]
together with their conjugates.
The modular indicator formula
\cite[Theorem~7.5]{NgSchauenburgExponent} gives respectively
\[
 \begin{split}
 \bigl(\nu_2(X_k)\bigr)_{k=0}^6&=(1,5,5,10,9,9,9),\\
 \bigl(\nu_2(X_k)\bigr)_{k=0}^6&=(1,1,1,2,1,1,1).
 \end{split}
\]
Every source simple is self-dual, so each second indicator must be
$1$ or $-1$.  Index $3$ contradicts this in all four cases.

The verifier \path{verify_braided_frontier.py}, in the continuation's
\path{criteria/} directory, binds the canonical tensors and
replays the signed sphericalization and positive Galois reduction,
fusion closures, character assignments, complete twist domain and
exact indicators.  Both normal and optimized executions pass.
Fibonacci and semion indicator controls give $(1,1)$ and $(1,-1)$,
so negative indicators are explicitly allowed.  This certificate
excludes braidings and leaves the general ledger unchanged.
\end{proof}
